\documentclass[11pt]{article}
\usepackage[margin=2.5cm]{geometry}
\usepackage{amsmath}
\usepackage[colorlinks,citecolor=blue,urlcolor=blue]{hyperref}
\usepackage{booktabs}
\usepackage{graphicx}
\usepackage[version=4]{mhchem}

\title{Simulations of Phytic Acid near Oleosin Covered Oil Drop\\
       \large Working document}
\author{}
\date{}

\begin{document}
\maketitle

\section{System}

We model the N-terminal arm of \textit{Cannabis sativa} oleosin
(UniProt A0A803P977) grafted to a planar surface that mimics the oil
body monolayer~\cite{Tzen1992}.  The FASTA sequence of the modelled
fragment is \texttt{MKSFLPESGPSAQQ}, where lowercase \texttt{n}
and \texttt{c} denote coarse-grained N- and C-terminal beads and
\texttt{a} is an anchor bead tethered to the surface.
The 14 amino acid residues are represented at one-bead-per-residue
resolution using the CALVADOS~3 force field~\cite{Tesei2021,Tesei2023}.
Consecutive beads are connected by harmonic bonds with force constant
$k = 80.33$\,kJ\,mol$^{-1}$\,\AA$^{-2}$ and equilibrium distance
$r_\text{eq} = 3.8$\,\AA.

$N = 50$ chains are grafted at surface density
$\sigma \approx 0.003$\,\AA$^{-2}$ (brush regime) on the bottom wall
of a slit geometry ($130 \times 130 \times 200$\,\AA$^{3}$, periodic
in $xy$, hard walls at $z = \pm 100$\,\AA).  Grafting is realised by
a harmonic tether on the anchor bead at $z_0 \approx -100$\,\AA\ with
$k = 0.5$\,kJ\,mol$^{-1}$\,\AA$^{-2}$.

\begin{figure}[h]
\centering
\includegraphics[width=0.7\textwidth]{grafted.png}
\caption{Simulation snapshot showing the oleosin brush (bottom) with a
single phytate molecule (red, top). Beads are colored by charge.}
\label{fig:snapshot}
\end{figure}

\section{Peptide titration}

Four residue types undergo pH-dependent protonation via atom-type swap
moves with implicit protons ($a_{\text{H}^+} = 10^{-\text{pH}}$).
The protonated and deprotonated forms differ only in charge:

\begin{table}[h]
\centering
\begin{tabular}{@{}lllc@{}}
\toprule
Residue & Protonated (charge) & Deprotonated (charge) & $\text{p}K_\text{a}$ \\
\midrule
GLU  & HGLU (0)   & GLU ($-1$)  & 4.24 \\
LYS  & LYS ($+1$) & LYS0 (0)   & 10.4 \\
NTR  & HNTR ($+1$) & NTR (0)   & 8.0  \\
CTR  & HCTR (0)   & CTR ($-1$)  & 3.67 \\
\bottomrule
\end{tabular}
\caption{Peptide residue titration reactions.
  $\text{p}K_\text{a}$ values from Ref.~\cite{Grimsley2009}.}
\end{table}

\begin{table}[h]
\centering
\begin{tabular}{@{}ccrc@{}}
\toprule
pH & $I$ (M) & $\langle Z \rangle$/chain & $\sigma_Z$ \\
\midrule
 4 & 0.03 & $+0.61$ & 0.64 \\
 4 & 0.15 & $+0.84$ & 0.68 \\
 7 & 0.03 & $-0.10$ & 0.30 \\
 7 & 0.15 & $-0.12$ & 0.33 \\
10 & 0.03 & $-1.05$ & 0.33 \\
10 & 0.15 & $-1.19$ & 0.42 \\
\bottomrule
\end{tabular}
\caption{Mean charge per peptide chain averaged over all umbrella windows.
  $\sigma_Z$ is the standard deviation of the per-chain charge distribution
  (charge capacitance$^{1/2}$).}
\end{table}

\section{Phytate titration}

One molecule of myo-inositol hexakisphosphate (phytate, IP6) is
included as a titrating ligand.  Each molecule is coarse-grained as
seven beads: one inositol centre (INO, $\sigma = 6.2$\,\AA,
$M = 180$\,Da) and six phosphate groups (PH, $\sigma = 5.8$\,\AA,
$M = 95$\,Da) placed at the crystallographic P positions in the
chair conformation ($\approx 4.5$\,\AA\ from the ring centre, one
axial and five equatorial).  Bead sizes are scaled from CALVADOS~3
by $M^{1/3}$ to maintain volume consistency with the amino acid beads.

Phytate has 12 dissociable protons (two per phosphate group),
giving 13 charge states from H$_{12}$PA (neutral) to PA$^{12-}$.
All 12 dissociation steps are modelled via molecular swap moves
in the semi-grand canonical ensemble~\cite{SmithTriska1994}.
The general dissociation reaction is
\[
\ce{H_{\textit{n}}PA^{(12-\textit{n})-} <=> H_{\textit{n}-1}PA^{(13-\textit{n})-} + H+}
\]
with macroscopic $\text{p}K_{\text{a},n}$ values from potentiometric
measurements in NaCl solution~\cite{DeStefano2003}:

\begin{table}[h]
\centering
\begin{tabular}{@{}cccccccccccccc@{}}
\toprule
$n$ & 12 & 11 & 10 & 9 & 8 & 7 & 6 & 5 & 4 & 3 & 2 & 1 \\
\midrule
$\text{p}K_{\text{a},n}$ & 1.1 & 1.5 & 1.7 & 2.1 & 2.5 & 2.9 & 5.72 & 6.28 & 6.81 & 7.60 & 9.94 & 11.84 \\
\bottomrule
\end{tabular}
\end{table}

\noindent
The net charge is distributed uniformly over the six phosphate beads;
protons are treated implicitly with activity $a_{\ce{H+}} =
10^{-\text{pH}}$.  All 13 protonation states share the same molecular
geometry (seven beads), enabling molecular swap moves that preserve
spatial orientation via gyration-tensor alignment.

\section{Interactions}

Nonbonded interactions between all beads combine the Ashbaugh--Hatch
potential~\cite{Ashbaugh2008} (short-range, hydrophobicity-dependent)
with screened Coulomb (Debye--H\"uckel) electrostatics.
For phytate--amino acid cross-pairs, the short-range potential is
replaced by WCA (purely repulsive) since phytate beads are
non-hydrophobic.

Intramolecular nonbonded pairs separated by one bond are excluded
(\texttt{excluded\_neighbours:~1}).  All intramolecular pairs within
each phytate molecule are excluded.
The aqueous medium is modelled with the temperature-dependent
dielectric permittivity of water at 298.15\,K.

\section{Monte Carlo moves}

All simulations were performed with a Rust reimplementation of
Faunus~\cite{Stenqvist2013} (\url{https://github.com/mlund/faunus-rs}).
Sampling uses the Metropolis criterion with the following moves per sweep:

\begin{itemize}
\item \textbf{Rigid-body translation} of peptide chains in $xy$
      ($\delta_\text{max} = 10.0$\,\AA) and phytate molecules along $z$
      ($\delta_\text{max} = 5.0$\,\AA, 20 repeats).
\item \textbf{Pivot rotation} of peptide chains
      ($\delta\theta_\text{max} = 0.5$\,rad).
\item \textbf{Rigid-body rotation} of phytate molecules
      ($\delta\theta_\text{max} = 1.0$\,rad).
\item \textbf{Speciation (molecular swap)} for the 12 phytate
      titration reactions and 4 peptide atom-type swaps
      (10 repeats per draw).
\end{itemize}

\section{Potential of mean force}

The potential of mean force (PMF) for phytate as a function of
distance from the brush surface is computed using multi-walker
umbrella sampling with hard-wall windows.  The collective variable
is the $z$-position of the INO bead, spanning the full slit
$[-L_z/2, +L_z/2]$ with 20\,\AA\ wide windows and 10\,\AA\ spacing
(50\% overlap).

Each window starts from a common equilibrated state.  A harmonic
drive phase ($k_\text{drive} = 5.0$\,kJ\,mol$^{-1}$\,\AA$^{-2}$)
steers the phytate into the target window; the bias is then removed
for unbiased production sampling.  Per-window states are saved for
restart.

Adjacent-window histograms are stitched via the overlap-ratio method:
\[
\Delta F_{i \to i+1} = -k_BT \ln \frac{f_i}{f_{i+1}}
\]
where $f_i$ is the fraction of samples in window $i$ that fall within
the overlap region with window $i{+}1$.

PMFs are computed at pH~4, 7, and 10 with ionic strengths
$I = 0.03$\,M, $0.15$\,M, and $1000$\,M NaCl.

\begin{figure}[h]
\centering
\includegraphics[width=\textwidth]{pmf_scan.png}
\caption{Potential of mean force for phytate as a function of distance
from the grafted oleosin brush.
Left: effect of pH at $I = 0.15$\,M.
Right: effect of ionic strength at pH~7.
The PMF is shifted so that the midplane region ($|z| < 20$\,\AA)
averages to zero.}
\label{fig:pmf}
\end{figure}

\bibliographystyle{unsrt}
\bibliography{references}
\end{document}
